\documentclass[../../main.tex]{subfiles}% 注意这里的文件路径不能用 ./main.tex ,否则用latexmk编译子文件会报错
\graphicspath{{\subfix{./image/}}{\subfix{../../image/}}} % 指定图片引用路径,后续可以直接使用图片文件名
% 如果图片存放在上级目录,则使用 ../../image/ 作为备用路径

\begin{document}

\section{与函数相伴的 Young 测度序列}\label{section:ch05-收敛}

设 $X$ 与 $Y$ 为具可数基的局部紧 Hausdorff 空间，$\mu\in\mathscr{R}^+_f(X)$，$u_j:X\to Y$ 为 $(\mathscr{B}(X),\mathscr{B}(Y))$-可测（$j\in\mathbb{N}$），$\nu_{u_j}\in\mathcal{Y}^+(X;Y)$ 为与 $u_j$ 相伴的 Young 测度（见\cref{definition:ch02-Young测度}）。把\cref{subsection:ch05-Young弱星}与\cref{subsection:ch05-Young窄}中关于 $\mathcal{Y}^+(X;Y)$ 的一般结果应用于序列 $\{\nu_{u_j}\}$，可获得关于序列 $\{u_j\}$ 本身收敛性的信息。以下沿用前文的记号：$\mathcal{C}(X\times Y)$、$\mathcal{C}_b(X\times Y)$ 与 $\mathcal{C}_0(X\times Y)$ 分别记 Carathéodory 函数、有界 Carathéodory 被积函数与 $C_0$-Carathéodory 被积函数所成的空间（见\cref{definition:ch05-Caratheodory}）。此方向的第一个结果如下。

\begin{proposition}\label{proposition:ch05-几乎处处收敛}\label{proposition:ch05-5.4.1}
设 $X$ 与 $Y$ 为具可数基的局部紧 Hausdorff 空间，$\mu\in\mathscr{R}^+_f(X)$。设 $u_j,u:X\to Y$ 为 $(\mathscr{B}(X),\mathscr{B}(Y))$-可测（$j\in\mathbb{N}$），$\nu_{u_j},\nu_u\in\mathcal{Y}^+(X;Y)$ 分别为与 $u_j$ 及 $u$ 相伴的 Young 测度。则：
\begin{enumerate}[(i)]
\item\label{item:ch05-ae-i} 若 $u_j\to u$ 在 $X$ 中 $\mu$-几乎处处成立，则 $\nu_{u_j}\stackrel{n}{\rightharpoonup}\nu_u$ 在 $\mathcal{Y}^+(X;Y)$ 中；
\item\label{item:ch05-ae-ii} 若 $\nu_{u_j}\stackrel{n}{\rightharpoonup}\nu_u$ 在 $\mathcal{Y}^+(X;Y)$ 中，则存在子列 $\{u_{j_k}\}\subseteq\{u_j\}$ 使 $u_{j_k}\to u$ 在 $X$ 中 $\mu$-几乎处处成立。
\end{enumerate}
\end{proposition}

\begin{proof}
先作准备：由\cref{proposition:appA-局部紧基}，具可数基的局部紧 Hausdorff 空间可度量化，故 $Y$ 上存在与拓扑相容的度量 $d$。

(i) 由\eqref{align:ch05-5.53}，对每个 $j\in\mathbb{N}$ 有 $(\nu_{u_j})_x=\delta_{u_j(x)}$ 对 $\mu$-几乎所有 $x\in X$，且 $(\nu_u)_x=\delta_{u(x)}$ 对 $\mu$-几乎所有 $x\in X$。由于 $\mu$ 有限，可数多个零测集之并仍为零测集，故存在 $\mu$-零测集 $N\subseteq X$ 使对所有 $x\in N^c$ 与所有 $j\in\mathbb{N}$ 同时成立上述两式。对这样的 $x$，由\eqref{align:ch05-5.15}得
\begin{align*}
d_P\big((\nu_{u_j})_x,(\nu_u)_x\big)=d_P\big(\delta_{u_j(x)},\delta_{u(x)}\big)=\min\{1,d(u_j(x),u(x))\},
\end{align*}
其中 $d_P$ 为 $\mathscr{R}^+_f(Y)$ 上的 Prokhorov 度量（见\cref{definition:ch05-Prokhorov}）。由假设 $u_j(x)\to u(x)$，故 $\min\{1,d(u_j(x),u(x))\}\to0$，即 $d_P\big((\nu_{u_j})_x,(\nu_u)_x\big)\to0$。于是由\cref{lemma:ch05-Prokhorov到窄}，对每个 $x\in N^c$ 有 $(\nu_{u_j})_x\stackrel{n}{\rightharpoonup}(\nu_u)_x$ 在 $\mathscr{R}^+_f(Y)$ 中。再由\cref{proposition:ch05-逐点窄收敛}得 $\nu_{u_j}\stackrel{n}{\rightharpoonup}\nu_u$ 在 $\mathcal{Y}^+(X;Y)$ 中。

(ii) 映射 $h:X\times Y\to\mathbb{R}$，$h(x,y)\triangleq\min\{1,d(y,u(x))\}$，属于 $\mathcal{C}_b(X\times Y)$：对每个 $x\in X$，$h(x,\cdot)=\min\{1,d(\cdot,u(x))\}$ 连续且取值于 $[0,1]$；对每个 $y\in Y$，映射 $x\mapsto d(y,u(x))$ 为 $\mathscr{B}(X)$-可测，因为它是 $(\mathscr{B}(X),\mathscr{B}(Y))$-可测映射 $u$ 与连续函数 $d(y,\cdot)$ 的复合；且 $\|h(x,\cdot)\|_\infty\leqslant1\in L^1(X)$。由假设 $\nu_{u_j}\stackrel{n}{\rightharpoonup}\nu_u$ 与\cref{definition:ch05-Young窄}，
\begin{align*}
\lim_{j\to\infty}\int_X\mathrm{d}\mu(x)\int_Y\min\{1,d(y,u(x))\}\,\mathrm{d}(\nu_{u_j})_x(y)
=\int_X\mathrm{d}\mu(x)\int_Y\min\{1,d(y,u(x))\}\,\mathrm{d}(\nu_u)_x(y),
\end{align*}
而由\eqref{align:ch05-5.54}（分别用于 $u_j$ 与 $u$）上式即
\begin{align*}
\lim_{j\to\infty}\int_X\min\{1,d(u_j(x),u(x))\}\,\mathrm{d}\mu(x)
=\int_X\min\{1,d(u(x),u(x))\}\,\mathrm{d}\mu(x)=0.
\end{align*}
被积函数非负且不超过 $1$，故由 Chebyshev 不等式，对每个 $\varepsilon\in(0,1)$，
\begin{align*}
\mu\big(\{\min\{1,d(u_j,u)\}\geqslant\varepsilon\}\big)
\leqslant\frac{1}{\varepsilon}\int_X\min\{1,d(u_j(x),u(x))\}\,\mathrm{d}\mu(x)\to0,
\end{align*}
即 $\min\{1,d(u_j,u)\}\to0$ 依测度。于是由\cref{proposition:ch02-收敛关系}(ii)，存在子列 $\{u_{j_k}\}\subseteq\{u_j\}$ 使 $\min\{1,d(u_{j_k}(x),u(x))\}\to0$ 对 $\mu$-几乎所有 $x\in X$。对这样的 $x$，当 $k$ 充分大时 $\min\{1,d(u_{j_k}(x),u(x))\}<1$，故它等于 $d(u_{j_k}(x),u(x))$；因此 $d(u_{j_k}(x),u(x))\to0$，即 $u_{j_k}(x)\to u(x)$。结论得证。
\end{proof}

% SUPPLEMENT: parent=proposition:ch05-几乎处处收敛; kind=motivation; 动机（几乎处处收敛与窄收敛的桥梁）、(i) 与 (ii) 关于子列的不对称性、集中现象与一致可积模、有限测度与可数基假设的作用、反例说明 (ii) 不能加强为全序列
% SUPPLEMENT: parent=proposition:ch05-5.4.1; kind=motivation; 别名 label（供后文引用其 (ii)），动机、条件作用与反例见下一条注记
\begin{remark}\phantomsection\label{remark:ch05-收敛注}
\cref{proposition:ch05-几乎处处收敛}把序列 $\{u_j\}$ 的几乎处处收敛与其 Young 测度序列的窄收敛联系起来，两个方向有本质差别。

(i) 方向 (i) 无需取子列：若 $u_j\to u$ 在 $X$ 中 $\mu$-几乎处处成立，则 $\nu_{u_j}\stackrel{n}{\rightharpoonup}\nu_u$。其含义是，窄收敛只检验有界 Carathéodory 被积函数 $h$ 的积分，而 $u_j\to u$ 几乎处处与 $h$ 关于 $y$ 的连续性结合即可用控制收敛定理取极限。

(ii) 方向 (ii) 只能对子列成立：窄收敛只控制积分 $\int_Xh(x,u_j(x))\,\mathrm{d}\mu(x)$，不控制 $u_j$ 的逐点值。证明表明，由 $\nu_{u_j}\stackrel{n}{\rightharpoonup}\nu_u$ 只能推出 $\min\{1,d(u_j,u)\}\to0$ 依测度，而依测度收敛只保证存在几乎处处收敛的子列（见\cref{proposition:ch02-收敛关系}(ii)）。

(iii) 集中现象。若极限 $\nu\in\mathcal{Y}^+(X;\mathbb{R})$ 不与任何 $\mathscr{B}(X)$-可测函数相伴（这样的 $\nu$ 存在，见\cref{proposition:ch05-非相伴}），则 $\{u_j\}$ 不可能既在 $L^1(X)$ 中有界又一致可积：此时其一致可积模 $\eta(\{u_j\})$ 为正（见\cref{definition:ch02-一致可积模}），质量沿 $u_j$ 的振幅方向集中，而窄极限把这份集中质量保留为 $\nu_x$ 的非原子部分。

(iv) 反例（(ii) 不能加强为全序列收敛）。取 $X=(0,1)$，$\mu=\lambda|_{(0,1)}$，$u\equiv0$，并置 $E_j\triangleq\big(\tfrac{k}{2^m},\tfrac{k+1}{2^m}\big)$，$u_j\triangleq\chi_{E_j}$，其中 $j=2^m+k$，$0\leqslant k<2^m$。则 $\mu(\{u_j\neq u\})=2^{-m}\to0$，故对每个 $h\in\mathcal{C}_b((0,1)\times\mathbb{R})$，
\begin{align*}
\left|\int_0^1h(x,u_j(x))\,\mathrm{d}x-\int_0^1h(x,0)\,\mathrm{d}x\right|\leqslant2\|h\|_\infty2^{-m}\to0,
\end{align*}
即 $\nu_{u_j}\stackrel{n}{\rightharpoonup}\nu_0=\delta_0$；然而对每个 $x\in(0,1)$，点 $x$ 属于无穷多个 $E_j$ 且不属于无穷多个 $E_j$，故 $\{u_j(x)\}$ 不收敛。因此取子列是必要的；在此例中 $j_k=2^k$ 给出 $u_{j_k}\to0$ 在 $(0,1)$ 中 $\mu$-几乎处处。

(v) 假设的作用与边界：$\mu$ 的有限性用于把可数多个例外零测集并为零测集，并保证 $\|h(x,\cdot)\|_\infty\in L^1(X)$；$X$ 与 $Y$ 的可数基假设保证 $Y$ 可度量化，从而 Prokhorov 度量 $d_P$ 与\eqref{align:ch05-5.15}可用，且 $\mathcal{Y}^+(X;Y)$ 上的窄收敛可由序列刻画。若无可数基假设，$Y$ 未必可度量化，$d_P$ 论证失效，须改以紧集上的一致结构替代。
\end{remark}

\subsection{$\{\nu_{u_j}\}$ 的弱$^*$收敛}\label{subsection:ch05-函数弱星}

以下处理 $Y=\mathbb{R}$ 的情形；对 $Y=\mathbb{R}^d$（$d\in\mathbb{N}$，$d\geqslant2$）类似结果成立。设 $F$ 为 $X$ 到 $\mathbb{R}$ 的 $\mathscr{B}(X)$-可测函数族，并置
\begin{align*}
N_F\triangleq\{\nu_u\mid u\in F\}\subseteq\mathcal{Y}^+(X;\mathbb{R}).
\end{align*}
由\cref{corollary:ch05-Young一一对应}与\cref{proposition:ch05-Dirac测度}，每个 $\nu_u\in N_F$ 满足 $\|(\nu_u)_x\|_{\mathscr{R}_f(\mathbb{R})}=1$ 对 $\mu$-几乎所有 $x\in X$，故 $N_F$ 含于空间 $L^\infty_{w^*}(X;\mathscr{R}_f(\mathbb{R}))$ 的单位球
\begin{align*}
U\triangleq\left\{\nu\in L^\infty_{w^*}(X;\mathscr{R}_f(\mathbb{R}))\ \middle|\
\operatorname*{ess\,sup}_{x\in X}\|\nu_x\|_{\mathscr{R}_f(\mathbb{R})}\leqslant1\right\}
\end{align*}
之中（见\eqref{align:ch05-5.55}）。此外，由\cref{theorem:ch04-可分Y对偶}（取 $Y=C_0(\mathbb{R})$，此时 $Y^*=\mathscr{R}_f(\mathbb{R})$，见\cref{theorem:ch02-Riesz表示}）与\cref{proposition:ch05-Caratheodory同构}，
\begin{align*}
L^\infty_{w^*}(X;\mathscr{R}_f(\mathbb{R}))=\big(L^1(X;C_0(\mathbb{R}))\big)^*=\big(\mathcal{C}_0(X\times\mathbb{R})\big)^*,
\end{align*}
且由\cref{theorem:ch04-Lp可分}，空间 $L^1(X;C_0(\mathbb{R}))$ 可分，故 Banach–Alaoglu 定理可用于 $U$。于是有如下结果（见\cite{ref08}）。

\begin{theorem}[Ball]\label{theorem:ch05-ball}
设 $X$ 为具可数基的局部紧 Hausdorff 空间，$\mu\in\mathscr{R}^+_f(X)$。设 $F$ 为 $X$ 到 $\mathbb{R}$ 的 $\mathscr{B}(X)$-可测函数族。则对任意序列 $\{u_j\}\subseteq F$：
\begin{enumerate}[(i)]
\item\label{item:ch05-ball-i} 存在子列 $\{u_{j_k}\}\subseteq\{u_j\}$ 与 $\nu\in U$（记 $\nu\equiv\{\nu_x\}_{x\in X}$）使
\begin{align}
g\circ u_{j_k}\stackrel{*}{\rightharpoonup}g_*\quad\text{在 }L^\infty(X)\text{ 中，对所有 }g\in C_0(\mathbb{R}),\label{align:ch05-5.66}
\end{align}
其中
\begin{align}
g_*(x)\triangleq\int_{\mathbb{R}}g(y)\,\mathrm{d}\nu_x(y)\quad\text{对 }\mu\text{-几乎所有 }x\in X;\label{align:ch05-5.67}
\end{align}
\item\label{item:ch05-ball-ii} 若对某个闭子集 $K\subseteq\mathbb{R}$ 有
\begin{align}
\lim_{j\to\infty}\mu\big(u_j^{-1}(A^c)\big)=0\quad\text{对 }K\text{ 的任意开邻域 }A,\label{align:ch05-5.68}
\end{align}
则 $\operatorname{supp}\nu_x\subseteq K$ 对 $\mu$-几乎所有 $x\in X$。
\end{enumerate}
\end{theorem}

\begin{proof}
(i) 由 Banach–Alaoglu 定理，对任意 $\{\nu_{u_j}\}\subseteq N_F\subseteq U$，存在 $\{\nu_{u_{j_k}}\}\subseteq\{\nu_{u_j}\}$ 与 $\nu\in U$ 使 $\nu_{u_{j_k}}\stackrel{*}{\rightharpoonup}\nu$ 在 $L^\infty_{w^*}(X;\mathscr{R}_f(\mathbb{R}))$ 中，即
\begin{align}
\lim_{k\to\infty}\int_X\mathrm{d}\mu(x)\int_{\mathbb{R}}h(x,y)\,\mathrm{d}(\nu_{u_{j_k}})_x(y)
=\int_X\mathrm{d}\mu(x)\int_{\mathbb{R}}h(x,y)\,\mathrm{d}\nu_x(y)\label{align:ch05-5.69}
\end{align}
对所有 $h\in L^1(X;C_0(\mathbb{R}))$ 成立。在\eqref{align:ch05-5.69}中取 $h(x,y)=f(x)g(y)$，其中 $f\in L^1(X)$，$g\in C_0(\mathbb{R})$，得
\begin{align*}
\lim_{k\to\infty}\int_X\mathrm{d}\mu(x)f(x)\int_{\mathbb{R}}g(y)\,\mathrm{d}(\nu_{u_{j_k}})_x(y)
=\int_X\mathrm{d}\mu(x)f(x)\int_{\mathbb{R}}g(y)\,\mathrm{d}\nu_x(y),
\end{align*}
故由\eqref{align:ch05-5.54}
\begin{align*}
\lim_{k\to\infty}\int_Xf(x)g(u_{j_k}(x))\,\mathrm{d}\mu(x)
=\int_X\mathrm{d}\mu(x)f(x)\int_{\mathbb{R}}g(y)\,\mathrm{d}\nu_x(y)
\end{align*}
对所有 $f\in L^1(X)$ 成立。由\eqref{align:ch05-5.49}，映射 $x\mapsto\int_{\mathbb{R}}g(y)\,\mathrm{d}\nu_x(y)$ 为 $\mathscr{B}(X)$-可测；此外
\begin{align*}
\operatorname*{ess\,sup}_{x\in X}\left|\int_{\mathbb{R}}g(y)\,\mathrm{d}\nu_x(y)\right|\leqslant\|g\|_\infty,
\end{align*}
故由\eqref{align:ch05-5.67}定义的函数 $g_*$ 属于 $L^\infty(X)$。断言得证。

(ii) 若 $K=\mathbb{R}$，则断言成立。设 $K\subset\mathbb{R}$ 满足\eqref{align:ch05-5.68}，$g\in C_0(\mathbb{R})$ 且 $g|_K=0$。对任意 $\varepsilon>0$，集合 $A_\varepsilon\triangleq\{y\in\mathbb{R}\mid|g(y)|<\varepsilon\}$ 为开集且 $A_\varepsilon\supset K$。于是由假设
\begin{align*}
\lim_{j\to\infty}\mu\big(u_j^{-1}((A_\varepsilon)^c)\big)
=\lim_{j\to\infty}\mu\big(\{x\in X\mid|g(u_j(x))|\geqslant\varepsilon\}\big)=0,
\end{align*}
故 $g(u_j)\to0$ 依测度。由于 $\{g(u_j)\}$ 在 $L^\infty(X)$ 中有界且 $\mu$ 有限，由控制收敛定理（见\cref{proposition:ch02-强收敛性质}(i)）得
\begin{align*}
\lim_{j\to\infty}\int_Xf(x)g(u_j(x))\,\mathrm{d}\mu(x)=0\quad\text{对所有 }f\in L^1(X).
\end{align*}
由该等式与\eqref{align:ch05-5.66}–\eqref{align:ch05-5.67}得 $\int_{\mathbb{R}}g(y)\,\mathrm{d}\nu_x(y)=0$ 对 $\mu$-几乎所有 $x\in X$。由于 $g|_K=0$，结论成立。
\end{proof}

% SUPPLEMENT: parent=theorem:ch05-ball; kind=proof; 原证明末句“由于 g|_K=0，结论成立”省略了从 ∫g dν_x=0（对所有在 K 上为零的 g∈C_0(ℝ)）到 supp ν_x ⊆ K 的推理；补出 Urysohn 截断、紧集覆盖与可数并的完整论证
\begin{proof}[补充证明：\cref{theorem:ch05-ball}(ii) 的最后一步]
由于 $\nu$ 是正测度序列 $\{\nu_{u_{j_k}}\}\subseteq\mathcal{Y}^+(X;\mathbb{R})$ 的弱$^*$极限，有 $\nu_x\in\mathscr{R}^+_f(\mathbb{R})$ 对 $\mu$-几乎所有 $x\in X$。已证：对每个满足 $g|_K=0$ 的 $g\in C_0(\mathbb{R})$，
\begin{align*}
\int_{\mathbb{R}}g(y)\,\mathrm{d}\nu_x(y)=0\quad\text{对 }\mu\text{-几乎所有 }x\in X.
\end{align*}
固定 $l\in\mathbb{N}$，置 $C_l\triangleq K^c\cap[-l,l]$。由于 $K$ 闭，$C_l$ 为紧集。对每个 $y\in C_l$，由\cref{proposition:appA-Urysohn}存在 $g_y\in C_c(\mathbb{R})$ 使 $0\leqslant g_y\leqslant1$，$g_y|_K=0$，$g_y(y)=1$；于是 $A_y\triangleq\{g_y>\tfrac12\}$ 为 $y$ 的开邻域。由 $C_l$ 的紧性，存在 $y_1,\cdots,y_m\in C_l$ 使 $C_l\subseteq\bigcup_{i=1}^mA_{y_i}$。对 $\mu$-几乎所有 $x\in X$，由 $g_{y_i}\geqslant\tfrac12$ 在 $A_{y_i}$ 上得 $\nu_x(A_{y_i})\leqslant2\int_{\mathbb{R}}g_{y_i}\,\mathrm{d}\nu_x=0$，于是
\begin{align*}
\nu_x(C_l)\leqslant\sum_{i=1}^m\nu_x(A_{y_i})=0.
\end{align*}
由于上述论证只涉及可数多个函数 $g_y$，对应的例外零测集之并仍为零测集；又 $K^c=\bigcup_{l=1}^\infty C_l$，故 $\nu_x(K^c)=0$ 对 $\mu$-几乎所有 $x\in X$。因此 $\operatorname{supp}\nu_x\subseteq K$ 对 $\mu$-几乎所有 $x\in X$。
\end{proof}

% SUPPLEMENT: parent=theorem:ch05-ball; kind=motivation; 动机（只要求可测性、不要求 L^1 有界）、极限落在单位球 U 与质量损失的含义、条件 (5.68) 的必要性反例、K=ℝ 的平凡情形、与 Thm 5.4.5 的分工
\begin{remark}\phantomsection\label{remark:ch05-ball注}
\cref{theorem:ch05-ball}是本分节的基本紧性结果：对 $X$ 到 $\mathbb{R}$ 的任意 $\mathscr{B}(X)$-可测函数列 $\{u_j\}$（不要求 $\{u_j\}$ 在 $L^1(X)$ 中有界，甚至不要求 $u_j$ 可积），相应的 Young 测度序列总有弱$^*$收敛的子列。代价是极限 $\nu$ 只属于单位球 $U$：可能发生质量损失，即 $\|\nu_x\|_{\mathscr{R}_f(\mathbb{R})}<1$，从而 $\nu$ 不必属于 $\mathcal{Y}^+(X;\mathbb{R})$。

关于 (i)：\eqref{align:ch05-5.66}断言对每个 $g\in C_0(\mathbb{R})$，复合函数 $g\circ u_{j_k}$ 在 $L^\infty(X)$ 中弱$^*$收敛于 $g_*(x)=\int_{\mathbb{R}}g\,\mathrm{d}\nu_x$，即所有在无穷远消失的非线性复合函数的弱$^*$极限都可由极限测度 $\nu$ 表示。取 $g\in C_c(\mathbb{R})$ 逼近 $\chi_F$（$F\in\mathscr{B}(\mathbb{R})$ 且 $\nu_x(\partial F)=0$）可读出 $\nu_x(F)$。

关于 (ii)：\eqref{align:ch05-5.68}是质量不沿 $y$ 方向逃向无穷的精确条件：它要求对 $K$ 的每个开邻域 $A$，$\{u_j\}$ 最终 $\mu$-几乎全部落入 $u_j^{-1}(A)$。若去掉该条件，结论不成立：取 $X=\{x_0\}$，$\mu=\delta_{x_0}$，$K=\{0\}$，$u_j\equiv1$，则 $\nu_{x_0}=\delta_1$，$\operatorname{supp}\nu_{x_0}=\{1\}\not\subseteq K$；而此时 $A=(-\tfrac12,\tfrac12)$ 是 $K$ 的开邻域且 $\mu(u_j^{-1}(A^c))=\mu(X)=1$ 不趋于 $0$，故\eqref{align:ch05-5.68}确实失效。

当 $K=\mathbb{R}$ 时\eqref{align:ch05-5.68}自动成立，结论 (ii) 平凡。若进一步假设 $\{u_j\}$ 在 $L^1(X)$ 中有界，则极限不损失质量，并可得到窄收敛的结论与\eqref{align:ch05-5.74}（见\cref{theorem:ch05-函数窄收敛}）；因此\cref{theorem:ch05-ball}与\cref{theorem:ch05-函数窄收敛}的分工是：前者只要可测性、极限在 $U$ 中，后者要 $L^1$ 有界、极限在 $\mathcal{Y}^+(X;\mathbb{R})$ 中。
\end{remark}

\subsection{$\{\nu_{u_j}\}$ 的窄收敛}\label{subsection:ch05-函数窄}

现将\cref{subsection:ch05-Young窄}的结果应用于与函数相伴的 Young 测度序列。此方向的一个预备步骤是如下结果。

\begin{proposition}\label{proposition:ch05-有界紧}\label{proposition:ch05-弱星紧}
设 $X$ 为具可数基的局部紧 Hausdorff 空间，$\mu\in\mathscr{R}^+_f(X)$。设 $F\subseteq L^1(X)$ 满足 $\sup_{u\in F}\|u\|_{L^1(X)}<\infty$（即 $F$ 在 $L^1(X)$ 中有界）。则与 $F$ 相伴的 Young 测度集合 $N_F\triangleq\{\nu_u\in\mathcal{Y}^+(X;\mathbb{R})\mid u\in F\}$ 在 $\mathcal{Y}^+(X;\mathbb{R})$ 中胎紧（见\cref{definition:ch05-Young胎紧}）。
\end{proposition}

\cref{proposition:ch05-有界紧}是如下引理的直接推论。

\begin{lemma}\label{lemma:ch05-胎紧等价}
设 $X$ 为具可数基的局部紧 Hausdorff 空间，$\mu\in\mathscr{R}^+_f(X)$，$F\subseteq L^1(X)$。则下列命题等价：
\begin{enumerate}[(i)]
\item\label{item:ch05-tightfun-i} 集合 $N_F\triangleq\{\nu_u\in\mathcal{Y}^+(X;\mathbb{R})\mid u\in F\}$ 在 $\mathcal{Y}^+(X;\mathbb{R})$ 中胎紧；
\item\label{item:ch05-tightfun-ii} 存在下紧函数 $\psi:[0,\infty)\to[0,\infty)$（见\cref{definition:ch05-下紧}）使
\begin{align}
\sup_{u\in F}\int_X\psi(|u(x)|)\,\mathrm{d}\mu(x)<\infty;\label{align:ch05-5.70}
\end{align}
\item\label{item:ch05-tightfun-iii} $\lim\limits_{k\to\infty}\Big(\sup\limits_{u\in F}\mu(\{|u|\geqslant k\})\Big)=0$，其中 $\{|u|\geqslant k\}\triangleq\{x\in X\mid|u(x)|\geqslant k\}$。
\end{enumerate}
\end{lemma}

\begin{proof}
(i)$\Rightarrow$(ii)。由\cref{proposition:ch05-Young胎紧判据}，存在下紧函数 $\varphi:\mathbb{R}\to[0,\infty)$ 使
\begin{align}
\sup_{\nu_u\in N_F}\int_{\mathbb{R}}\varphi(y)\,\mathrm{d}(\nu_u)_{p_{\mathbb{R}}}(y)<\infty,\label{align:ch05-5.71}
\end{align}
其中 $(\nu_u)_{p_{\mathbb{R}}}$ 为 $\nu_u$ 在 $\mathbb{R}$ 上的投影（见\eqref{align:ch05-5.33}）。由\cref{remark:ch05-径向下紧}，在\cref{proposition:ch05-Young胎紧判据}的证明中可取 $\varphi(y)=\psi(|y|)$，其中 $\psi(t)\triangleq\sum_{n=0}^\infty n\,\chi_{(n,n+1]}(t)$（$t\in[0,\infty)$）。函数 $\psi$ 下紧，且由\eqref{align:ch05-5.71}
\begin{align}
\sup_{\nu_u\in N_F}\int_{\mathbb{R}}\psi(|y|)\,\mathrm{d}(\nu_u)_{p_{\mathbb{R}}}(y)<\infty.\label{align:ch05-5.72}
\end{align}
另一方面，由\eqref{align:ch05-5.53}与\eqref{align:ch05-5.36}，对任意 $u\in F$，
\begin{align}
\int_X\psi(|u(x)|)\,\mathrm{d}\mu(x)
&=\int_X\mathrm{d}\mu(x)\int_{\mathbb{R}}\psi(|y|)\,\mathrm{d}(\nu_u)_x(y)\notag\\
&=\int_{X\times\mathbb{R}}\psi(|y|)\,\mathrm{d}\nu_u(x,y)
=\int_{\mathbb{R}}\psi(|y|)\,\mathrm{d}(\nu_u)_{p_{\mathbb{R}}}(y).\label{align:ch05-5.73}
\end{align}
由\eqref{align:ch05-5.72}–\eqref{align:ch05-5.73}得\eqref{align:ch05-5.70}，断言得证。

(ii)$\Rightarrow$(i)。设 $\psi:[0,\infty)\to[0,\infty)$ 下紧。则映射 $\varphi:\mathbb{R}\to[0,\infty)$，$\varphi(y)\triangleq\psi(|y|)$，亦下紧：对每个 $r\in\mathbb{R}$，集合 $S_r\triangleq\{t\in[0,\infty)\mid\psi(t)\leqslant r\}$ 紧，而 $\{\varphi\leqslant r\}=\{y\in\mathbb{R}\mid|y|\in S_r\}=S_r\cup(-S_r)$ 为紧集。于是由\eqref{align:ch05-5.70}与\eqref{align:ch05-5.73}得\eqref{align:ch05-5.72}，从而得\eqref{align:ch05-5.71}。由\cref{proposition:ch05-Young胎紧判据}，断言得证。

(ii)$\Rightarrow$(iii)。不妨设 $\psi$ 非减：否则以 $\psi_*(t)\triangleq\inf_{s\geqslant t}\psi(s)$ 代替 $\psi$，它非减、满足 $\psi_*\leqslant\psi$、下紧，且由 $\psi$ 的下紧性得 $\psi(t)\to\infty$（$t\to\infty$），故\eqref{align:ch05-5.70}对 $\psi_*$ 仍成立。由 Chebyshev 不等式，对任意 $k>0$ 与 $u\in F$，
\begin{align*}
\mu(\{|u|\geqslant k\})\leqslant\frac{1}{\psi(k)}\int_{\{|u|\geqslant k\}}\psi(|u(x)|)\,\mathrm{d}\mu(x)
\leqslant\frac{1}{\psi(k)}\int_X\psi(|u(x)|)\,\mathrm{d}\mu(x),
\end{align*}
于是由\eqref{align:ch05-5.70}，对每个 $u\in F$ 有
\begin{align*}
\mu(\{|u|\geqslant k\})\leqslant\frac{1}{\psi(k)}\,\sup_{v\in F}\int_X\psi(|v(x)|)\,\mathrm{d}\mu(x);
\end{align*}
由于 $\psi(k)\to\infty$，令 $k\to\infty$ 得 (iii)。

(iii)$\Rightarrow$(ii)。由假设存在非减发散序列 $\{k_l\}\subseteq[0,\infty)$ 使
\begin{align*}
\sup_{u\in F}\mu(\{|u|\geqslant k_l\})\leqslant\frac{1}{l^3}\quad\text{对所有 }l\in\mathbb{N}.
\end{align*}
置 $k_0\triangleq0$ 与 $\psi(t)\triangleq\sum_{l=0}^\infty l\,\chi_{(k_l,k_{l+1}]}(t)$（$t\in[0,\infty)$）。对 $r\geqslant0$，取 $n\triangleq[r]$，则 $t>k_{n+1}$ 蕴含 $\psi(t)\geqslant n+1>r$，故 $\{\psi\leqslant r\}\subseteq[0,k_{n+1}]$ 为紧集，即 $\psi$ 下紧。此外，对所有 $u\in F$，
\begin{align*}
\int_X\psi(|u(x)|)\,\mathrm{d}\mu(x)
=\sum_{l=0}^\infty l\,\mu(\{k_l\leqslant|u|<k_{l+1}\})
\leqslant\sum_{l=1}^\infty l\,\mu(\{|u|\geqslant k_l\})
\leqslant\sum_{l=1}^\infty\frac{1}{l^2}<\infty.
\end{align*}
证明完成。
\end{proof}

% SUPPLEMENT: parent=lemma:ch05-胎紧等价; kind=proof; 原证明 (i)=>(ii) 直接断言“由 Remark 5.1.10 可取 φ(y)=ψ(|y|)”而未说明该径向函数的来源；此处补出由任一下紧 φ 构造非减下紧径向权并保持 (5.71) 的严格论证
\begin{proof}[补充证明：\cref{lemma:ch05-胎紧等价}(i)$\Rightarrow$(ii) 中径向权函数的构造]
设 $\varphi:\mathbb{R}\to[0,\infty)$ 下紧且满足\eqref{align:ch05-5.71}。置 $\tilde\psi(t)\triangleq\inf_{|y|\geqslant t}\varphi(y)$（$t\in[0,\infty)$）。则 $\tilde\psi$ 非负、非减，且 $\tilde\psi(|y|)\leqslant\varphi(y)$ 对每个 $y\in\mathbb{R}$ 成立。由于 $\varphi$ 下紧，对每个 $r\geqslant0$ 存在 $\rho_r>0$ 使 $\{\varphi\leqslant r\}\subseteq[-\rho_r,\rho_r]$；于是 $t>\rho_r$ 蕴含 $\tilde\psi(t)>r$，即 $\{\tilde\psi\leqslant r\}\subseteq[0,\rho_r]$ 为紧集，故 $\tilde\psi$ 下紧，且 $\tilde\psi(t)\to\infty$（$t\to\infty$）。由 $\tilde\psi(|y|)\leqslant\varphi(y)$ 得
\begin{align*}
\sup_{\nu_u\in N_F}\int_{\mathbb{R}}\tilde\psi(|y|)\,\mathrm{d}(\nu_u)_{p_{\mathbb{R}}}(y)
\leqslant\sup_{\nu_u\in N_F}\int_{\mathbb{R}}\varphi(y)\,\mathrm{d}(\nu_u)_{p_{\mathbb{R}}}(y)<\infty,
\end{align*}
即\eqref{align:ch05-5.72}对 $\tilde\psi$ 成立。以 $\tilde\psi$ 代替阶梯函数 $\psi$，\eqref{align:ch05-5.73}与\eqref{align:ch05-5.70}的推导不变。
\end{proof}

\begin{proof}[证明\cref{proposition:ch05-有界紧}]
由假设 $\sup_{u\in F}\|u\|_{L^1(X)}=\sup_{u\in F}\int_X|u|\,\mathrm{d}\mu<\infty$，故在\eqref{align:ch05-5.70}中可取 $\psi:[0,\infty)\to[0,\infty)$，$\psi(y)=y$；该函数下紧，因为 $\{\psi\leqslant r\}=[0,r]$ 对每个 $r\geqslant0$ 为紧集。于是由\cref{lemma:ch05-胎紧等价}(ii)$\Rightarrow$(i)，$N_F$ 胎紧。
\end{proof}

\begin{remark}\phantomsection\label{remark:ch05-vallee比较}
\cref{lemma:ch05-胎紧等价}(ii) 比 de la Vallée-Poussin 判别法（见\cref{proposition:ch02-有界一致可积等价}(iii)）弱。这与\cref{proposition:ch02-有界一致可积等价}和\cref{proposition:ch05-有界紧}一致：$F$ 的有界性蕴含\cref{lemma:ch05-胎紧等价}(ii)（取 $\psi(y)=y$），而 de la Vallée-Poussin 判别法等价于 $F$ 的有界性与一致可积性二者。
\end{remark}

% SUPPLEMENT: parent=proposition:ch05-有界紧; kind=motivation; 动机（把 L^1 有界性翻译为胎紧性，从而接上定理 5.3.11）、假设的作用、与定理 5.4.2 的分工、去掉有界性的反例
% SUPPLEMENT: parent=proposition:ch05-弱星紧; kind=motivation; 别名 label（沿用旧译名），动机与各项条件的作用见下一条注记
% SUPPLEMENT: parent=lemma:ch05-胎紧等价; kind=motivation; 动机（三条判据的含义与相互翻译）、非退化条件、反例说明 (iii) 严格弱于一致可积
\begin{remark}\phantomsection\label{remark:ch05-胎紧判据函数}
\cref{proposition:ch05-有界紧}与\cref{lemma:ch05-胎紧等价}的作用是把函数序列的紧性问题化为测度序列的紧性问题。

(i) \cref{proposition:ch05-有界紧}的假设只有一条：$F$ 在 $L^1(X)$ 中有界；结论 $N_F$ 胎紧正是\cref{theorem:ch05-弱紧Young}的假设，故二者结合给出窄拓扑中的相对序列紧性（\cref{theorem:ch05-函数窄收敛}）。若去掉有界性，结论可为假：取 $X=(0,1)$，$\mu=\lambda|_{(0,1)}$，$u_j\equiv j$，则 $\nu_{u_j}=\mu\otimes\delta_j$ 不胎紧（对任意紧集 $K\subseteq\mathbb{R}$ 与 $j\notin K$ 有 $\nu_{u_j}((0,1)\times K^c)=1$），其弱$^*$极限为 $0$（见\cref{remark:ch05-ball注}）。

(ii) \cref{lemma:ch05-胎紧等价}的三条判据分别刻画：胎紧性 (i) 即 $\mathcal{Y}^+(X;\mathbb{R})$ 中的胎紧（\cref{definition:ch05-Young胎紧}）；(ii) 是“存在在无穷远增长的下紧权 $\psi$，使 $\psi(|u|)$ 的积分关于 $u\in F$ 一致有界”，即矩条件；(iii) 是“尾部集合 $\{|u|\geqslant k\}$ 的测度关于 $u\in F$ 一致地趋于 $0$”。证明给出三者之间的显式转换，权函数形如 $\psi(t)=\sum_nn\chi_{(n,n+1]}(t)$ 或 $\psi(t)=\sum_ll\chi_{(k_l,k_{l+1}]}(t)$。

(iii) 条件 (iii) 严格弱于一致可积性：取 $X=(0,1)$，$\mu=\lambda|_{(0,1)}$，$u_j\triangleq j\chi_{(0,1/j)}$（$j\in\mathbb{N}$）。则 $\|u_j\|_{L^1(X)}=1$；对 $k>0$，当 $j\geqslant k$ 时 $\mu(\{|u_j|\geqslant k\})=1/j$，当 $j<k$ 时该测度为 $0$，故 $\sup_{j}\mu(\{|u_j|\geqslant k\})=1/k\to0$，即 (iii) 成立。然而 $\int_{\{|u_j|\geqslant t\}}|u_j|\,\mathrm{d}\mu=1$ 对一切 $j\geqslant t$ 成立，故 $\{u_j\}$ 不一致可积。这与\cref{remark:ch05-vallee比较}一致：胎紧性只对应 de la Vallée-Poussin 判别法中较弱的一侧。

(iv) 非退化条件与边界：$\psi$ 只要求下紧，不要求连续、凸或满足 $\psi(t)/t\to\infty$；(iii) 中的极限自动为 $0$，无需对 $\{u_j\}$ 附加逐点假设。$\mu$ 的有限性不可去：若 $\mu(X)=\infty$，则\eqref{align:ch05-5.70}与 (iii) 的等价性失效，须先对 $\mu$ 作 $\sigma$-有限的局部化。
\end{remark}

由\cref{proposition:ch05-有界紧}与\cref{theorem:ch05-弱紧Young}，若 $F\subseteq L^1(X)$ 有界，则相应的 Young 测度集合 $N_F$ 在窄拓扑中相对序列紧。于是有如下结果。

\begin{theorem}\label{theorem:ch05-函数窄收敛}
设 $X$ 为具可数基的局部紧 Hausdorff 空间，$\mu\in\mathscr{R}^+_f(X)$。设 $F\subseteq L^1(X)$ 满足 $\sup_{u\in F}\|u\|_{L^1(X)}<\infty$（即 $F$ 在 $L^1(X)$ 中有界）。则对任意序列 $\{u_j\}\subseteq F$，存在子列 $\{u_{j_k}\}\subseteq\{u_j\}$ 与 Young 测度 $\nu\in\mathcal{Y}^+(X;\mathbb{R})$ 使对所有有界 Carathéodory 被积函数 $h\in\mathcal{C}_b(X\times\mathbb{R})$ 成立如下积分等式：
\begin{align}
\lim_{k\to\infty}\int_Xh(x,u_{j_k}(x))\,\mathrm{d}\mu(x)
=\int_X\mathrm{d}\mu(x)\int_{\mathbb{R}}h(x,y)\,\mathrm{d}\nu_x(y).\label{align:ch05-5.74}
\end{align}
\end{theorem}

\begin{proof}
固定任意 $\{u_j\}\subseteq F$。由\cref{proposition:ch05-有界紧}，相应的 Young 测度集合 $N_F$ 胎紧，故由\cref{theorem:ch05-弱紧Young}存在子列 $\{\nu_{u_{j_k}}\}\subseteq\{\nu_{u_j}\}$ 与 Young 测度 $\nu\in\mathcal{Y}^+(X;\mathbb{R})$ 使 $\nu_{u_{j_k}}\stackrel{n}{\rightharpoonup}\nu$，即对所有 $h\in\mathcal{C}_b(X\times\mathbb{R})$，
\begin{align*}
\lim_{k\to\infty}\int_{X\times\mathbb{R}}h(x,y)\,\mathrm{d}\nu_{u_{j_k}}(x,y)
=\int_{X\times\mathbb{R}}h(x,y)\,\mathrm{d}\nu(x,y).
\end{align*}
由\eqref{align:ch05-5.54}，
\begin{align*}
\int_{X\times\mathbb{R}}h(x,y)\,\mathrm{d}\nu_{u_{j_k}}(x,y)
=\int_Xh(x,u_{j_k}(x))\,\mathrm{d}\mu(x),
\end{align*}
而由\eqref{align:ch05-5.36}，
\begin{align*}
\int_{X\times\mathbb{R}}h(x,y)\,\mathrm{d}\nu(x,y)
=\int_X\mathrm{d}\mu(x)\int_{\mathbb{R}}h(x,y)\,\mathrm{d}\nu_x(y).
\end{align*}
故等式\eqref{align:ch05-5.74}成立。
\end{proof}

% SUPPLEMENT: parent=theorem:ch05-函数窄收敛; kind=motivation; 动机（有界性保证极限不损失质量）、假设的作用、(5.74) 的含义、与定理 5.4.2 的区别、极限不必与函数相伴
\begin{remark}\phantomsection\label{remark:ch05-窄紧性}
\cref{theorem:ch05-函数窄收敛}是\cref{theorem:ch05-ball}在 $L^1$ 有界情形下的加强：$F$ 的 $L^1$ 有界性经\cref{proposition:ch05-有界紧}给出 $N_F$ 的胎紧性，再经\cref{theorem:ch05-弱紧Young}给出窄收敛的子列，因此极限不损失质量，即 $\nu\in\mathcal{Y}^+(X;\mathbb{R})$ 而非仅 $\nu\in U$（比较\cref{theorem:ch05-ball}）。

关于假设：$\sup_{u\in F}\|u\|_{L^1(X)}<\infty$ 只用于胎紧性；若去掉它而保留 $u_j$ 的可测性，则只得到\cref{theorem:ch05-ball}的弱$^*$结论，极限可能为 $0$。$\mu$ 的有限性用于胎紧判据（\cref{lemma:ch05-胎紧等价}）。

关于\eqref{align:ch05-5.74}：该式说明 $\nu$ 与子列 $\{u_{j_k}\}$ 相伴，即子列的所有有界 Carathéodory 复合函数的积分极限都由 $\nu$ 的离散分解给出，这与\cref{proposition:ch05-Dirac测度}中 $\nu_u$ 的形式一致。但 $\nu$ 不必与任何 $\mathscr{B}(X)$-可测函数相伴（见\cref{proposition:ch05-非相伴}与\cref{example:ch05-周期Young测度}），故\eqref{align:ch05-5.74}不能反过来给出 $u_{j_k}$ 的几乎处处收敛。

边界情形：若 $\{u_j\}$ 在 $L^1(X)$ 中有界且弱收敛，则\eqref{align:ch05-5.74}的右端仍由 $\nu$ 给出，但 $u_{j_k}$ 不必强收敛；此时重心 $u_*(x)=\int_{\mathbb{R}}y\,\mathrm{d}\nu_x(y)$ 是弱极限的候选（见\cref{remark:ch05-biting注}）。
\end{remark}

\cref{theorem:ch05-函数窄收敛}的如下推论只需非负的 Carathéodory 被积函数。

\begin{corollary}\label{corollary:ch05-下半连续}
设 $X$ 为具可数基的局部紧 Hausdorff 空间，$\mu\in\mathscr{R}^+_f(X)$。设 $F\subseteq L^1(X)$ 满足 $\sup_{u\in F}\|u\|_{L^1(X)}<\infty$。则对任意 $\{u_j\}\subseteq F$，存在子列 $\{u_{j_k}\}\subseteq\{u_j\}$ 与 Young 测度 $\nu\in\mathcal{Y}^+(X;\mathbb{R})$ 使对任意 Carathéodory 函数 $h:X\times\mathbb{R}\to[0,\infty)$（见\cref{definition:ch05-Caratheodory}(i)），
\begin{align}
\int_{X\times\mathbb{R}}h(x,y)\,\mathrm{d}\nu(x,y)
\leqslant\liminf_{k\to\infty}\int_Xh(x,u_{j_k}(x))\,\mathrm{d}\mu(x).\label{align:ch05-5.75}
\end{align}
\end{corollary}

\begin{proof}
对任意 $\{u_j\}\subseteq F$，由\cref{theorem:ch05-函数窄收敛}存在子列 $\{\nu_{u_{j_k}}\}\subseteq\{\nu_{u_j}\}$ 与 Young 测度 $\nu\in\mathcal{Y}^+(X;\mathbb{R})$ 使 $\nu_{u_{j_k}}\stackrel{n}{\rightharpoonup}\nu$。于是由\eqref{align:ch05-5.62}（用于非负 Carathéodory 函数 $h$）得
\begin{align*}
\int_{X\times\mathbb{R}}h(x,y)\,\mathrm{d}\nu(x,y)
\leqslant\liminf_{k\to\infty}\int_X\mathrm{d}\mu(x)\int_{\mathbb{R}}h(x,y)\,\mathrm{d}(\nu_{u_{j_k}})_x(y),
\end{align*}
而由\eqref{align:ch05-5.54}，上式右端的内部积分等于 $h(x,u_{j_k}(x))$。故\eqref{align:ch05-5.75}成立。
\end{proof}

% SUPPLEMENT: parent=corollary:ch05-下半连续; kind=motivation; 动机（无界非负 Carathéodory 检验函数的下半连续性）、非负性假设的必要性、与 Thm 5.4.5 的分工、变号时的反例线索
\begin{remark}\phantomsection\label{remark:ch05-下半连续估计}
\cref{corollary:ch05-下半连续}是\cref{proposition:ch05-Caratheodory下半连续}与\cref{theorem:ch05-函数窄收敛}的结合：它把窄收敛的检验函数类从有界的 $\mathcal{C}_b(X\times\mathbb{R})$ 扩大到无界但非负的 Carathéodory 函数，代价是把等式减弱为\eqref{align:ch05-5.75}中的不等式。

关于假设：$h\geqslant0$ 不可去。\eqref{align:ch05-5.62}的证明只对非负 $h$ 成立（它对 $h$ 作截断 $h_n=\min\{\inf_\zeta[h+n\,d(\zeta,\cdot)],n\}$ 并用单调收敛定理），若 $h$ 变号，$\int h\,\mathrm{d}\nu$ 可能严格大于 $\liminf$：取 $X=(0,1)$，$\mu=\lambda$，$u_j\triangleq j\chi_{(0,1/j)}$，$h(y)\triangleq-y$，则 $u_j\to0$ 依测度，$\nu_x=\delta_0$，故 $\int_{X\times\mathbb{R}}h\,\mathrm{d}\nu=0$，而 $\int_Xh(x,u_j(x))\,\mathrm{d}x=-1$ 对一切 $j$，故 $0=\int_{X\times\mathbb{R}}h\,\mathrm{d}\nu>-1=\liminf_{j\to\infty}\int_Xh(x,u_j(x))\,\mathrm{d}x$，即\eqref{align:ch05-5.75}对变号的 $h$ 失效。

关于结论的形式：\eqref{align:ch05-5.75}中的 $\liminf$ 不能换为极限（同上例，取 $h(y)=|y|$ 时 $\int_Xh(x,u_j(x))\,\mathrm{d}x\equiv1$ 而 $\int h\,\mathrm{d}\nu=0$，极限存在但与左端不等）。该不等式在下一小节用于去掉\cref{corollary:ch05-下半连续}中 $h$ 的非负性限制（见\cref{proposition:ch05-5.4.10}）。
\end{remark}

把\cref{theorem:ch05-函数窄收敛}（或\cref{theorem:ch05-ball}，因为窄收敛蕴含弱$^*$收敛）给出的极限测度 $\nu\in\mathcal{Y}^+(X;\mathbb{R})$ 视为与该定理所述子列 $\{u_{j_k}\}$ 相伴是自然的。有时可以显式算出这一测度，如下结果所示。

\begin{proposition}\label{proposition:ch05-周期平均}
设 $I\subseteq\mathbb{R}$ 为长度为 $1$ 的区间，$K\subseteq\mathbb{R}$ 为非空紧集，$u:\mathbb{R}\to K$ 为 $\mathscr{B}(\mathbb{R})$-可测且以 $1$ 为周期。置 $u_k(x)\triangleq u(kx)$（$k\in\mathbb{N}$，$x\in I$），$\{\nu_{u_k}\}\subseteq\mathcal{Y}^+(I;K)$ 为相应的 Young 测度序列，其中测度空间取为 $(I,\mathscr{B}(I),\lambda|_I)$。则 $\nu_{u_k}$ 在 $\mathcal{Y}^+(I;K)$ 中窄收敛于乘积测度 $\nu=\lambda|_I\times\lambda_u$，其中 $\lambda_u\triangleq u_*(\lambda|_I)$ 为 $\lambda|_I$ 在 $u$ 下的像测度，即
\begin{align*}
\int_Ig\circ u\,\mathrm{d}\lambda=\int_{\mathbb{R}}g\,\mathrm{d}\lambda_u
\quad\text{对所有 }g\in L^1(\mathbb{R},\mathscr{B}(\mathbb{R}),\lambda_u).
\end{align*}
\end{proposition}

\begin{proof}
为简便，设 $I=(0,1)$。由于 $K\subseteq\mathbb{R}$ 紧，$C_b(K)$ 可分，且由\cref{remark:ch05-紧空间同构}，$\mathcal{C}_b(I\times K)=L^1(I;C_b(K))$。于是由\cref{proposition:ch04-简单函数稠密}(ii)（取 $D\subseteq C_b(K)$ 为可数稠密子集），只需证明对所有 $f\in C_c(I)$ 与 $g\in C_b(K)$ 有
\begin{align}
\lim_{k\to\infty}\int_If(x)g(u_k(x))\,\mathrm{d}\lambda(x)
=\int_If(x)\,\mathrm{d}\lambda(x)\int_Ig(u(x))\,\mathrm{d}\lambda(x).\label{align:ch05-5.76}
\end{align}
固定 $k\in\mathbb{N}$，设 $\{I_1,\cdots,I_k\}$ 为 $I$ 的分割，其中 $I_m$ 为以 $\frac{m-1}{k}$ 与 $\frac mk$ 为端点的区间（$m=1,\cdots,k$）。作变量替换 $x=\frac{\xi+m-1}{k}$（$\xi\in I$）并用 $u$ 的 $1$-周期性得
\begin{align}
\int_If(x)g(u_k(x))\,\mathrm{d}\lambda(x)
&=\int_If(x)g(u(kx))\,\mathrm{d}\lambda(x)\notag\\
&=\frac1k\sum_{m=1}^k\int_If\Big(\frac{\xi+m-1}{k}\Big)g(u(\xi))\,\mathrm{d}\lambda(\xi).\label{align:ch05-5.77}
\end{align}
由于对每个 $\xi\in I$，
\begin{align*}
\lim_{k\to\infty}\frac1k\sum_{m=1}^kf\Big(\frac{\xi+m-1}{k}\Big)=\int_If(x)\,\mathrm{d}\lambda(x),
\end{align*}
且上式左端的绝对值不超过 $\|f\|_\infty$，由控制收敛定理（对 $\xi$ 积分）在\eqref{align:ch05-5.77}中令 $k\to\infty$ 得\eqref{align:ch05-5.76}。结论得证。
\end{proof}

% SUPPLEMENT: parent=proposition:ch05-周期平均; kind=motivation; 动机（周期震荡的显式极限）、假设的作用（周期 1、I 长为 1、K 紧、u 只可测）、边界（非周期或区间长度不为 1 时极限依赖于 x）、与 Example 5.4.5 的联系
\begin{remark}\phantomsection\label{remark:ch05-周期}
\cref{proposition:ch05-周期平均}说明周期震荡序列 $\{u(kx)\}$ 的 Young 测度极限与 $x$ 无关，等于“一个周期上的分布”$\lambda_u$ 与“位置测度”$\lambda|_I$ 的乘积。这正是周期 homogenization 中“局部极限是周期分布的混合”这一现象的概率表述。

关于假设：$I$ 的长度为 $1$ 保证 $\lambda_u$ 是概率测度，且分割 $\{I_m\}$ 的每个小区间经 $x\mapsto kx$ 恰好覆盖一个完整周期；$u$ 的 $1$-周期性保证 $u(kx)$ 在每个 $I_m$ 上取遍 $u$ 在一个周期上的全部值；$K$ 紧（从而 $C_b(K)$ 可分）保证可用\cref{remark:ch05-紧空间同构}把 $\mathcal{C}_b(I\times K)$ 与 $L^1(I;C_b(K))$ 等同，并使用\cref{proposition:ch04-简单函数稠密}(ii) 作乘积型检验。$u$ 只要求 $\mathscr{B}(\mathbb{R})$-可测，无需连续。

关于结论：$\nu=\lambda|_I\times\lambda_u$ 的离散分解为 $\nu_x=\lambda_u$ 对 $\lambda$-几乎每个 $x\in I$，它与 $x$ 无关；若 $u$ 非常数，则 $\lambda_u$ 非原子（当 $u$ 的分布函数连续时），此时 $\nu$ 不与任何 $\mathscr{B}(I)$-可测函数相伴，见\cref{example:ch05-周期Young测度}。

边界情形：若 $I$ 的长度 $L\neq1$，则极限为 $\lambda|_I\times\lambda_{u,L}$，其中 $\lambda_{u,L}$ 是 $\lambda|_{(0,L)}$ 在 $u$ 下的像，一般不再是概率测度；若 $u$ 不周期，则极限一般依赖于 $x$。
\end{remark}

如下例子给出\cref{proposition:ch05-周期平均}的两个显式计算。

\begin{example}\label{example:ch05-周期Young测度}
(i) 设 $I=(-\tfrac12,\tfrac12)$，$K=[-1,1]$，$u(x)=\sin(\pi x)$（$x\in I$）。由 Riemann 积分的初等结果，
\begin{align*}
\int_I(g\circ u)(x)\,\mathrm{d}\lambda(x)=\int_{-\frac12}^{\frac12}g(u(x))\,\mathrm{d}x
=\frac{1}{\pi}\int_{-1}^1\frac{g(y)}{\sqrt{1-y^2}}\,\mathrm{d}y.
\end{align*}
故对 $\lambda$-几乎每个 $x\in I$ 与所有 $F\in\mathscr{B}(\mathbb{R})\cap[-1,1]$，
\begin{align*}
\nu_x(F)=\frac{1}{\pi}\int_F\frac{\mathrm{d}y}{\sqrt{1-y^2}}=\text{常数}.
\end{align*}

(ii) 设 $I=(0,1)$，$K=\{-1,1\}$，$u(x)=\operatorname{sgn}[\sin(2\pi x)]$（$x\in I$）。则
\begin{align*}
\int_I(g\circ u)(x)\,\mathrm{d}\lambda(x)=\int_0^{\frac12}g(u(x))\,\mathrm{d}x+\int_{\frac12}^1g(u(x))\,\mathrm{d}x,
\end{align*}
故 $\lambda_u=\frac{\delta_1+\delta_{-1}}{2}$。于是对 $\lambda$-几乎每个 $x\in I$ 有
\begin{align*}
\nu_x=\frac12(\delta_1+\delta_{-1})=\text{常数}.
\end{align*}
\end{example}

须指出：由\cref{theorem:ch05-函数窄收敛}或\cref{theorem:ch05-ball}给出的极限测度 $\nu\in\mathcal{Y}^+(X;\mathbb{R})$ 不必与任何 $\mathscr{B}(X)$-可测函数相伴（\cref{example:ch05-周期Young测度}(ii) 即如此）。下述命题说明这一点，其证明留到下一小节。

\begin{proposition}\label{proposition:ch05-非相伴}
设 $X$ 为具可数基的局部紧 Hausdorff 空间，$\mu\in\mathscr{R}^+_f(X)$。设 $\{u_j\}\subseteq L^1(X)$，$u\in L^1(X)$ 满足 $u_j\rightharpoonup u$ 在 $L^1(X)$ 中但 $u_j\not\to u$ 在 $L^1(X)$ 中强收敛。设 $\{\nu_{u_j}\}$ 为相应的 Young 测度序列。则存在子列 $\{\nu_{u_{j_k}}\}\subseteq\{\nu_{u_j}\}$ 与 Young 测度 $\nu\in\mathcal{Y}^+(X;\mathbb{R})$ 使：
\begin{enumerate}[(i)]
\item\label{item:ch05-nonassoc-i} $\nu_{u_{j_k}}\stackrel{n}{\rightharpoonup}\nu$；
\item\label{item:ch05-nonassoc-ii} $\nu$ 不与任何 $\mathscr{B}(X)$-可测函数相伴。
\end{enumerate}
\end{proposition}

% SUPPLEMENT: parent=proposition:ch05-非相伴; kind=motivation; 动机（窄极限一般不是 Dirac 型）、假设 u_j⇀u 但非强收敛的作用、与 Example 5.4.5(ii) 的关系、边界（强收敛时必为 Dirac 型）
\begin{remark}\phantomsection\label{remark:ch05-非相伴注}
\cref{proposition:ch05-非相伴}说明\cref{theorem:ch05-函数窄收敛}的极限一般只是 Young 测度，不能写成 $\nu_f$ 的形式：弱收敛而非强收敛的序列会把“振荡”记录为极限测度的非原子部分。

关于假设：$u_j\rightharpoonup u$ 且 $u_j\not\to u$ 强收敛是本质的。若 $u_j\to u$ 强收敛于 $L^1(X)$，则由\cref{proposition:ch05-几乎处处收敛}(i)（沿几乎处处收敛的子列）与极限唯一性得 $\nu=\nu_u$，即 $\nu$ 与 $u$ 相伴，结论 (ii) 不成立。若去掉 $\{u_j\}$ 的 $L^1$ 有界性，则\cref{theorem:ch05-ball}的极限还可能为 $0$，此时 $\nu$ 同样不与任何函数相伴，但结论 (i) 中的窄收敛不再成立。

与\cref{example:ch05-周期Young测度}(ii) 的关系：该例中 $u$ 取两值、$\nu_x=\frac12(\delta_1+\delta_{-1})$ 非 Dirac 型，正是\cref{proposition:ch05-非相伴}的典型情形；而\cref{proposition:ch05-非相伴}指出这一现象在“弱收敛但非强收敛”的假设下是普遍存在的。

适用边界：结论 (ii) 中的“不相伴”是指不存在 $\mathscr{B}(X)$-可测函数 $f$ 使 $\nu=\nu_f$；由\cref{proposition:ch05-Dirac测度}，这等价于 $\nu_x$ 对 $\mu$-正测度集上的 $x$ 不是 Dirac 测度。
\end{remark}

\end{document}
